% ============================================================
%  Chapter 1 — Component Description
% ============================================================
\chapter{Component Description}
\label{ch:description}

\section{Provided Functions}
\label{sec:functions}

\textbf{resilient-call} is a Python library that provides two fault-tolerance
mechanisms as first-class building blocks:

\begin{itemize}[leftmargin=*, label=\textbullet]
    \item \textbf{Retry Policy} — automatically re-executes a failing callable
          up to a configurable number of times, with three backoff strategies:
          \textit{fixed}, \textit{exponential}, and \textit{jitter}.

    \item \textbf{Circuit Breaker} — monitors failure rates and, when a
          threshold is exceeded, \emph{opens} the circuit and short-circuits
          all subsequent calls without touching the downstream service.
          After a configurable timeout the breaker transitions to a probe state
          (\textsc{half-open}) and allows a test call through.
\end{itemize}

Both mechanisms can be used independently or combined through the
\texttt{ResilientExecutor} class, which composes them in the correct
execution order and optionally invokes a \textbf{fallback} callable when
all protective layers are exhausted.

The library additionally provides:

\begin{itemize}[leftmargin=*, label=\textbullet]
    \item Decorator API (\texttt{@retry}, \texttt{@circuit\_breaker}) for
          annotating existing functions with zero structural changes.
    \item A thread-safe named registry (\texttt{CircuitBreakerRegistry}) for
          sharing circuit breaker instances across modules.
    \item Full \texttt{asyncio} support via \texttt{execute\_async} methods.
    \item Event hooks: \texttt{on\_retry} and \texttt{on\_state\_change}
          callbacks for logging and monitoring integration.
\end{itemize}

\section{Purpose}
\label{sec:purpose}

Modern software systems are rarely self-contained — they communicate with
external services, databases, message brokers, and third-party APIs.  Any of
these dependencies can experience transient failures (network blips, temporary
overloads, rolling deployments) or complete outages.  Without explicit handling,
a single failing downstream service can propagate errors upward through the
entire call chain and bring down an otherwise healthy application.

\textbf{resilient-call} addresses this problem at the integration boundary.
It wraps outbound calls with configurable resilience logic so that callers
receive either a successful result (after automatic retries) or a safe default
value (fallback), while the downstream system is given time to recover.

\section{Business Problems Solved}
\label{sec:business-problems}

\begin{enumerate}[leftmargin=*, label=\arabic*.]
    \item \textbf{Service downtime causing revenue loss.}
          When a payments API is intermittently unavailable, unprotected code
          either throws an exception to the end user or retries aggressively,
          amplifying load on the failing service.  The circuit breaker pattern
          prevents both failure modes: it stops exposing failures to end users
          (fallback) and stops hammering the recovering service (OPEN state).

    \item \textbf{Cascading failures in microservice architectures.}
          A slow upstream service can exhaust the thread pool of all downstream
          callers.  By short-circuiting calls when a service is unhealthy, the
          circuit breaker prevents thread pool starvation from propagating.

    \item \textbf{Thundering herd after outages.}
          When a service recovers, all waiting clients retry simultaneously,
          potentially causing a second outage.  The JITTER backoff strategy and
          HALF-OPEN probe state both mitigate this risk.

    \item \textbf{Developer effort for boilerplate resilience code.}
          Each team writing its own retry loops and failure counters is
          duplicated work and a source of subtle bugs.  This component
          centralises the logic behind a well-tested, declarative API.
\end{enumerate}

\section{Intended Use}
\label{sec:intended-use}

\textbf{resilient-call} is intended for:

\begin{itemize}[leftmargin=*, label=\textbullet]
    \item Python backend services and microservices that call external HTTP
          APIs, gRPC services, message brokers, or databases.
    \item Any callable that may fail transiently and benefits from
          automatic retrying.
    \item Applications that need graceful degradation when a dependency is
          unavailable.
\end{itemize}

The primary usage patterns are shown in Listing~\ref{lst:basic-usage}.

\begin{lstlisting}[language=Python, caption={Basic usage — decorator and programmatic API}, label={lst:basic-usage}]
from resilient_call import (
    retry, circuit_breaker, BackoffStrategy,
    ResilientExecutor, RetryPolicy, CircuitBreaker,
)

# --- Decorator API ---

@retry(max_attempts=3, delay=0.5, backoff=BackoffStrategy.EXPONENTIAL)
def fetch_data(url: str) -> dict:
    return requests.get(url, timeout=5).json()

@circuit_breaker("payments-api", failure_threshold=5, timeout=30.0)
def charge_card(amount: float) -> dict:
    return payments_client.charge(amount)

# --- Programmatic API ---

executor = ResilientExecutor(
    retry_policy=RetryPolicy(max_attempts=4, delay=0.5),
    circuit_breaker=CircuitBreaker("inventory"),
    fallback=lambda exc: [],
)
items = executor.execute(fetch_inventory, warehouse_id=7)
\end{lstlisting}

\section{Restrictions}
\label{sec:restrictions}

\subsection{Technical Restrictions}

\begin{itemize}[leftmargin=*, label=\textbullet]
    \item \textbf{Python version:} Requires Python 3.10 or higher.
    \item \textbf{In-process state only:} Circuit breaker state is stored in
          process memory. In multi-replica deployments (e.g.\ Kubernetes pods)
          each replica maintains independent state.  For shared state, a
          distributed backing store (e.g.\ Redis) is required and is not
          provided by this component.
    \item \textbf{Thread-safety scope:} The component is safe for use from
          multiple threads within a single process, but is not optimised for
          very high call rates (above $\sim$100\,000 calls/s per breaker) due
          to lock contention.
    \item \textbf{Async support:} Full \texttt{asyncio} compatibility is
          provided, but mixing synchronous and asynchronous code paths via the
          same executor instance may cause event-loop conflicts in edge cases.
\end{itemize}

\subsection{Business Restrictions}

\begin{itemize}[leftmargin=*, label=\textbullet]
    \item The component does not provide authentication, authorisation, or
          encryption.  It operates purely at the call-reliability layer.
    \item Fallback functions are the caller's responsibility.  The component
          does not prescribe what the fallback should return.
    \item Retry logic is not suitable as a substitute for idempotency
          guarantees.  Callers must ensure that the wrapped operation is safe
          to retry (idempotent) or that duplicate executions are acceptable.
\end{itemize}

\section{Other Important Information}
\label{sec:other-info}

\begin{itemize}[leftmargin=*, label=\textbullet]
    \item \textbf{License:} MIT — free for commercial and non-commercial use.
    \item \textbf{Dependencies:} None.  The library uses only the Python
          standard library (\texttt{threading}, \texttt{time}, \texttt{random},
          \texttt{asyncio}, \texttt{dataclasses}, \texttt{enum}).
    \item \textbf{Distribution:} Installable via \texttt{pip install
          resilient-call}.
    \item \textbf{Source code:} Available on GitHub at
          \url{https://github.com/ardaerdogan/resilient-call}.
\end{itemize}
